% Kata pengantar. Tempat dan tanggal di bawah teks diambil dari
% config/identitas.tex (\kota dan \tanggalpengesahan).
\begin{katapengantar}
\panduan{Kata pengantar memuat pernyataan selesainya skripsi dan ucapan
terima kasih kepada pihak yang berperan langsung, misalnya dekan, kaprodi,
dosen pembimbing, dewan penguji, mitra penelitian, keluarga, dan rekan sejawat.
Tulis secara ringkas, khidmat, dan tanpa kalimat basa-basi berlebihan.}

Puji syukur dipanjatkan ke hadirat Tuhan Yang Maha Esa karena atas berkat dan rahmat-Nya,
naskah skripsi berjudul \enquote{\judulskripsi} ini dapat diselesaikan dengan baik sebagai
salah satu syarat memperoleh gelar Sarjana Komputer pada Program Studi Sistem Informasi,
Fakultas Ilmu Komputer, Universitas Darwan Ali.

Penyusunan naskah skripsi ini tidak terlepas dari bimbingan, dorongan, dan bantuan dari
berbagai pihak. Oleh karena itu, ucapan terima kasih dan penghargaan yang tulus disampaikan kepada:
\begin{enumerate}
  \item Dekan Fakultas Ilmu Komputer Universitas Darwan Ali;
  \item Ketua Program Studi Sistem Informasi Fakultas Ilmu Komputer;
  \item \dosendpembimbing\ selaku Dosen Pembimbing yang telah meluangkan waktu, tenaga, dan pikiran untuk memberikan bimbingan, arahan, serta dorongan berharga dalam penyelesaian skripsi ini;
  \item Ketua dan segenap Anggota Dewan Penguji atas saran, kritik konstruktif, dan masukan penyempurnaan naskah skripsi;
  \item Pimpinan, staf, dan seluruh jajaran organisasi mitra yang telah memberikan izin, fasilitas, serta kemudahan selama pelaksanaan penelitian dan akuisisi data;
  \item Orang tua, keluarga tercinta, dan seluruh kerabat yang senantiasa melimpahkan doa yang tulus, dorongan moral, cinta kasih, dan dukungan materiil yang tiada ternilai;
  \item Rekan-rekan mahasiswa Program Studi Sistem Informasi serta semua pihak yang tidak dapat disebutkan satu per satu, yang telah memberikan bantuan, semangat, dan kerja sama selama masa perkuliahan hingga penyusunan skripsi ini.
\end{enumerate}

Penulis menyadari bahwa naskah ini masih memiliki keterbatasan. Kritik dan saran yang membangun
senantiasa diharapkan demi penyempurnaan naskah ini. Semoga karya ini dapat memberikan manfaat
nyata bagi pengembangan keilmuan sistem informasi serta masyarakat luas.
\end{katapengantar}
